\chapter{Áp Lực Đồng Trang}
\markboth{Áp Lực Đồng Trang}{Áp Lực Đồng Trang}

\begin{truyen}{Tại Bạn Rủ Nên Em Mới Theo}{Peer Pressure}
\chuhoa{M}{ột bạn bị ghi tên vì quay clip chọc bạn trong lớp, liền chống chế: ``Dạ tại cả nhóm rủ em, chứ bình thường em không vậy. Em sợ bị tách ra nên em mới làm theo thôi.''}

\teacherpair{Thầy nhìn thẳng: ``Bạn bè có thể rủ. Nhưng tay em bấm nút quay là tay em. Miệng em cười phụ họa là miệng em. Áp lực đồng trang là có thật, nhưng nó không xóa phần lựa chọn của em.''}{The teacher says: ``Friends may invite you into it. But your hand still pressed record. Your mouth still laughed along. Peer pressure is real, but it does not erase your share of choice.''}

Bạn ấy nhỏ giọng: ``Nếu em không theo thì em bị nói là không cùng team.''

\teacherpair{Không cùng team trong một trò sai đôi khi lại là dấu hiệu đầu tiên cho thấy em đang đủ bản lĩnh để tự đứng bằng chân mình. Ai bắt em tệ đi để được công nhận thì không phải bạn, chỉ là đám đông đang cần thêm một người hùa.}{``Not being on the team in a wrong act is sometimes the first sign that you are strong enough to stand on your own. Anyone who needs you to become worse just to accept you is not a friend, only a crowd looking for one more follower.''}

\textit{(A student blames the crowd. The teacher acknowledges peer pressure but refuses to erase personal agency.)}
\end{truyen}

\begin{baihoc}
Đổ lỗi cho bạn bè là cách nhanh nhất để mình không bao giờ học được cách chịu trách nhiệm.
\textit{(Blaming friends is the fastest way to avoid learning responsibility.)}
\end{baihoc}

\begin{ghinhoanh}
\textbf{Short English to remember:}

Pressure is real.

Choice is still yours.
\end{ghinhoanh}

\begin{tuhocnhanh}
1. Em từng làm điều sai chỉ để không bị lạc nhóm chưa? \textit{(Have you ever done something wrong just to avoid being left out?)}

2. Em cần giữ nguyên tắc nào dù bạn bè không thích? \textit{(What principle should you keep even if your friends dislike it?)}
\end{tuhocnhanh}

\ngancach